\subsection{Financial Markets and Trading}

Financial markets move capital between those who hold it and those who need it. They allow participants to trade assets, hedge risk and raise funds, and each type of market serves a different purpose. Table~\ref{Tables:FinancialMarkets} summarises the main ones. The markets are connected, so a movement in one is often transmitted to the others and none of them is analysed in isolation. Prices are set by the participants themselves, from private investors to banks and investment funds, while regulators set the rules under which those participants operate. Electronic platforms and algorithmic execution have widened access and increased the speed at which prices adjust \cite{brandimarte_introduction_2018}.

\begin{table}[htb!]
\centering
\footnotesize
\begin{tabularx}{\textwidth}{@{}lXl@{}}
\toprule
\textbf{Market} & \textbf{Economic Role} & \textbf{Asset Examples} \\
\midrule
Stock Market & Capital raising through company shares trading. & Stocks, Equities \\
\addlinespace
Bond Market & Debt financing by entities borrowing from investors. & Government and Corporate Bonds \\
\addlinespace
Commodities Market & Trading of essential raw materials and agricultural goods. & Oil, Gold, Wheat \\
\addlinespace
Forex Market & Currency exchange for international trade. & USD, EUR, JPY, GBP \\
\addlinespace
Derivatives Market & Trading contracts for hedging or speculation. & Futures, Options, Swaps \\
\addlinespace
Indices Market & Aggregating and tracking market or sector performance. & S\&P 500, NASDAQ \\
\addlinespace
Cryptocurrency Market & Exchanging digital currencies. & Bitcoin, Ethereum \\
\bottomrule
\end{tabularx}
\caption{Summary of Financial Markets and respective Economic Role \cite{brandimarte_introduction_2018}.}
\label{Tables:FinancialMarkets}
\end{table}

\subsubsection{Forex Market: Characteristics and Significance}

The Foreign Exchange market, or Forex market, is where currencies are exchanged. It is the largest financial market by turnover, averaging 9.6 trillion US dollars per day in April 2025, a 28\% increase on the 7.5 trillion recorded three years earlier, and it operates continuously across time zones because the banks and institutions that quote its prices are spread worldwide. Spot transactions, the instrument traded in this work, account for 3 trillion of that daily total. The US Dollar sits on one side of 89.2\% of all trades, and the ten most traded currency pairs all involve it \cite{bis_otc_2025}.

Prices respond to interest rate decisions, economic releases and political events. Participants normally trade with leverage, which allows a position much larger than the deposited capital and scales profits and losses in the same proportion. Depth, continuous operation, sensitivity to news and the routine use of leverage are what make the market suitable for automated trading and hard to model at the same time. Its behaviour is noisy and non-stationary, which is why adaptive and data-driven strategies are of interest \cite{donnelly_art_2019}.

\subsubsection{Forex Trading: Fundamental Concepts and Operation}

Trading currency pairs involves buying one currency and simultaneously selling another. Buying the EUR/USD pair means purchasing Euros and selling US Dollars, in the expectation that the Euro will strengthen against the US Dollar. This work uses the seven major currency pairs: EUR/USD, USD/JPY, GBP/USD, AUD/USD, USD/CHF, NZD/USD and USD/CAD. They are built from the European Euro (EUR), US Dollar (USD), Japanese Yen (JPY), Great Britain Pound (GBP), Australian Dollar (AUD), Swiss Franc (CHF) and New Zealand Dollar (NZD). Every one of them has the US Dollar on one side, and deciding whether to buy or sell a pair amounts to predicting the strength of one currency relative to the other \cite{donnelly_art_2019}.

A position is the exposure held in a pair. A long position profits when the pair rises and loses when it falls; a short position does the opposite. Because a currency pair is a ratio between two currencies, selling is as natural as buying and requires no borrowing, which is not the case in equity markets. Positions are opened and closed with orders. A market order executes immediately at the price currently quoted. A limit order executes only at a specified price or better, and a stop order executes once the price passes a specified level. Two stop orders are commonly attached to an open position: a stop loss, which closes it once losses reach a chosen level, and a take profit, which closes it once a target is reached.

Forex trading occurs in a decentralised market that operates globally. There is no exchange where all transactions take place. Instead, a network of market makers, primarily banks (known as the interbank market) and brokers, facilitates trades. Interbank transactions involve large volumes and are dominated by institutions such as multinational corporations and major hedge funds. Retail traders participate through brokers with access to that network, who act as intermediaries. Market makers and brokers quote two prices at all times: the bid price, at which they are willing to buy the pair, and the ask price, at which they are willing to sell it. The difference between them is the spread. The absence of a central venue might suggest a lack of regulation, but competition among market makers keeps quoted prices close to one another and spreads narrow \cite{donnelly_art_2019}.

Traders read price through charts, and the candlestick chart is the standard form. Each chart is made up of candlesticks, which are structures that represent price movements over a time span that can range from minutes to years. A tick is a single update to the quoted price, published whenever a market maker revises its bid or its ask. Each candlestick, named for its resemblance to a double-sided candle, provides information on the opening price (O), the highest price (H), the lowest price (L), and the closing price (C) for a specific period of time (see Figure~\ref{Figure:CandlestickChart}). The colour of the candlestick represents the price movement, with lighter shades indicating an increase and darker shades indicating a decrease. This allows traders to gain a fast understanding of the direction of the market just by looking at the chart. Candlestick charts do not show volume (V), which is normally plotted separately below the price. Because Forex has no central venue, no consolidated record of traded size exists, so brokers report tick volume instead: the number of quote updates in the period. That is the measure used throughout this work \cite{rhoads_candlestick_2022}.

\begin{figure}[htb!]
\centering
\includegraphics[width=0.4\textwidth]{Images/CandlestickChart.png}
\caption{Candlesticks encoding upward (bullish) and downward (bearish) movements respectively \cite{teo_candlestick_nodate}.}
\label{Figure:CandlestickChart}
\end{figure}

This work uses tick data as its primary record, since it is the finest resolution the broker provides, and derives the opening, highest, lowest and closing prices together with the tick volume, or OHLCV data for short, whenever a coarser view is required \cite{rhoads_candlestick_2022}. Indicators computed from that data are used to give the model additional context, as described next.

\subsubsection{Forex Costs: Spreads, Commissions and Overnight Financing}

Trading is not free, and the charges are large relative to the profit available on a single position. Three of them apply in Forex, and they are charged independently of one another.

Prices are quoted in units that follow from the pair. A pip is the fourth decimal place of the quote for most pairs, or the second for pairs quoted against the Japanese Yen. A point is one tenth of a pip and is the smallest increment a broker quotes. Position size is expressed in units of the base currency, where one standard lot is $100{,}000$ units. For a position of volume $q$ at price $p$, the notional value in the quote currency is

\begin{equation}
V = q \cdot p
\end{equation}

The \textbf{spread} is the first charge and is paid implicitly. A position is opened at the ask and closed at the bid, so the round trip crosses the spread once. For a position of volume $q$, the cost is

\begin{equation}
C_{spread} = (p_{ask} - p_{bid}) \cdot q
\label{Equation:SpreadCost}
\end{equation}

The spread is not constant. It widens when liquidity falls, around economic releases and at the daily rollover, so a strategy tested with a fixed spread is tested against a market that does not exist.

The \textbf{commission} is the second charge and is explicit. Accounts that quote raw interbank spreads recover their margin through a commission charged on the traded volume rather than through a wider spread. Expressed as a number of points $c$, with $s_{point}$ the size of one point, the cost of a position of volume $q$ is

\begin{equation}
C_{commission} = c \cdot s_{point} \cdot q
\label{Equation:CommissionCost}
\end{equation}

The spread and the commission are separate charges and both are paid. Treating the spread as if it already included the commission understates the cost of trading on a raw-spread account.

The \textbf{swap}, or overnight financing, is the third charge and reflects the interest rate differential between the two currencies of the pair. It is applied when a position is still open at the daily rollover, and is charged at a multiple on one day of the week to account for the weekend. Depending on the direction of the position and the sign of the differential, it can be a cost or a credit. Some account types, offered so that clients can avoid interest for religious reasons, remove it entirely; these are known as swap-free accounts.

Finally, leverage determines how much capital a position requires. With leverage $L$, the margin committed to a position of notional value $V$ is

\begin{equation}
M = \frac{V}{L}
\end{equation}

Leverage does not change the profit or loss of a position, only the capital reserved against it, so it acts on the return of the account rather than on the return of the trade.

\subsubsection{Financial Analysis: Main Types and Formulas}

Representing the data is one half of the problem; the other is analysing it to decide what to trade. Financial analysis is generally divided into three types \cite{lim_handbook_2015}:

\begin{itemize}
\item \textbf{Fundamental Analysis}: examines economic factors that affect currency prices. This includes analysis of GDP growth, employment rates, interest rates, and monetary policies, which help predict currency trends based on a nation's economic health. However, the qualitative nature of these data poses challenges to computational models that prefer structured numerical input.
\item \textbf{Sentiment Analysis}: seeks to quantify the mood of the market by analysing news, market commentary, and similar qualitative data. It aims to measure how collective emotions impact market trends, although its complexity and need for advanced natural language processing make it less straightforward for model integration.
\item \textbf{Technical Analysis}: involves analysing previous OHLCV data to predict future market movements. It operates on the principle that past market behaviour will repeat itself, and thus identifying past patterns or metrics on current data can predict future movements. It is preferred for its quantitative results, which can be used by complex machine learning models.
\end{itemize}

Of the three, this work uses technical analysis, because it produces numerical series directly from price data and therefore needs no additional interpretation before a model can consume it. Technical analysis divides into indicators and candlestick patterns. Indicators are metrics created mainly by mathematicians and economists, calculated from historical market data, that offer a quantitative view of market dynamics. They are generally categorised into four main types \cite{lim_handbook_2015}:

\begin{itemize}
\item \textbf{Trend Indicators}: Used to identify the prevailing direction of the market by filtering out market noise. Refer to Table~\ref{Tables:TrendIndicators} for more details on frequently used indicators.
\item \textbf{Momentum Indicators}: Used to capture the momentum of a given trend, whether it is strong or weak, and if a reversal is expected. Refer to Table~\ref{Tables:MomentumIndicators} for more details on frequently used indicators.
\item \textbf{Volume Indicators}: Used to provide a perspective on the trading volume, providing clues about the strength of a price movement. Refer to Table~\ref{Tables:VolumeIndicators} for more details on frequently used indicators.
\item \textbf{Volatility Indicators}: Used to identify the level of uncertainty or risk associated with the volatility of the price movement. Refer to Table~\ref{Tables:VolatilityIndicators} for more details on frequently used indicators.
\end{itemize}

\begin{table}[htb!]
\centering
\footnotesize
\begin{tabularx}{\textwidth}{@{}lXl@{}}
\toprule
\textbf{Indicator Name} & \textbf{Formula} & \textbf{Range} \\
\midrule
Simple Moving Average & $\text{SMA}(N)_t = \frac{1}{N} \sum_{i=1}^{N} \text{Price}_{t-N+i}$ & - \\
\addlinespace
Exponential Moving Average & $\text{EMA}(N)_t = \frac{2}{N+1} \times \text{Price}_t + (1 - \frac{2}{N+1}) \times \text{EMA}(N)_{t-1}$ & - \\
\addlinespace
Weighted Moving Average & $\text{WMA}(N)_t = \frac{2}{N \times (N + 1)} \sum_{i=1}^{N} i \times \text{Price}_{t-N+i}$ & - \\
\addlinespace
Double Exponential Moving Average & $\text{DEMA}(N)_t = 2 \times \text{EMA}(N)_t - \text{EMA}_2(N)_t$ & - \\
\addlinespace
Triple Exponential Moving Average & $\text{TEMA}(N)_t = 3 \times [\text{EMA}(N)_t - \text{EMA}_2(N)_t] + \text{EMA}_3(N)_t$ & - \\
\addlinespace
Hull Moving Average & $\text{HMA}(N)_t = \text{WMA}\left(\sqrt{N}\right)\left[2 \times \text{WMA}\left(\frac{N}{2}\right)_t - \text{WMA}(N)_t\right]$ & - \\
\addlinespace
Triangular Moving Average & $\text{TRIMA}(N)_t = \frac{1}{N} \sum_{i=1}^{N} \text{SMA}(N)_{t-N+i}$ & - \\
\addlinespace
Efficiency Ratio & $\text{ER}(N)_t = \frac{| \text{Price}_t - \text{Price}_{t-N} |}{\sum_{i=1}^{N} | \text{Price}_{t-i+1} - \text{Price}_{t-i} |}$ & [0, 1] \\
\addlinespace
Kaufman Adaptive Moving Average & $\text{KAMA}(N)_t = \text{KAMA}(N)_{t-1} + \text{SC} \times (\text{Price}_t - \text{KAMA}(N)_{t-1})$ & - \\
\addlinespace
& $\text{SC} = [ \text{ER}(N)_t \times (\phi_{\text{short}} - \phi_{\text{long}}) + \phi_{\text{long}} ]^2, \phi_{\text{short}} = \frac{2}{N_{\text{short}} + 1}$, $\phi_{\text{long}} = \frac{2}{N_{\text{long}} + 1}$ & \\
\addlinespace
& where typically $N_{\text{short}} = 2$ and $N_{\text{long}} = 30$ & \\
\bottomrule
\end{tabularx}
\caption{Summary of frequently used Trend Technical Indicators \cite{jansen_machine_2020}.}
\label{Tables:TrendIndicators}
\end{table}
\begin{table}[htb!]
\centering
\footnotesize
\begin{tabularx}{\textwidth}{@{}lXl@{}}
\toprule
\textbf{Indicator Name} & \textbf{Formula} & \textbf{Range} \\
\midrule
Price Momentum & $\text{PM}_t = \frac{\text{Price}_t - \text{Price}_{t-1}}{\text{Price}_{t-1}}$ & - \\
\addlinespace
Rate of Change & $\text{ROC}(N)_t = \ln\left(\frac{\text{Price}_t}{\text{Price}_{t-N}}\right)$ & - \\
\addlinespace
Price Acceleration & $\text{PA}_t = \frac{\text{Short Linear Regression Coefficient}}{\text{Long Linear Regression Coefficient}} = \frac{\beta_\text{short}}{\beta_\text{long}}$ & - \\
\addlinespace
Relative Strength Index & $\text{RSI}(N)_t = 100 - \frac{100}{1 + \frac{\text{Average Gain}_t}{\text{Average Loss}_t}}$ & [0, 100] \\
\addlinespace
& $\text{Average Gain}_t = \frac{1}{N}\sum_{i=1}^{N} \max(\text{Price}_i - \text{Price}_{i-1}, 0)$ & \\
\addlinespace
& $\text{Average Loss}_t = \frac{1}{N}\sum_{i=1}^{N} \max(\text{Price}_{i-1} - \text{Price}_i, 0)$ & \\
\addlinespace
Aroon Oscillator & $\text{Aroon}(N)_t = \text{Aroon Up}(N)_t - \text{Aroon Down}(N)_t$ & [-100, 100] \\
\addlinespace
& $\text{Aroon Up}(N)_t = \left( \frac{N - \text{Periods Since Highest High in } N}{N} \right) \times 100$ & \\
\addlinespace
& $\text{Aroon Down}(N)_t = \left( \frac{N - \text{Periods Since Lowest Low in } N}{N} \right) \times 100$ & \\
\addlinespace
Balance of Power & $\text{BOP}_t = \frac{\text{Close}_t - \text{Open}_t}{\text{High}_t - \text{Low}_t}$ & [-1, 1] \\
\addlinespace
Commodity Channel Index & $\text{CCI}(N)_t = \frac{\text{TP}_t - \text{SMA}_\text{TP}(N)_t}{0.015 \times \text{Mean Deviation}(N)_t}$ & - \\
\addlinespace
& $\text{TP}_t = \frac{\text{High}_t + \text{Low}_t + \text{Close}_t}{3}$ & \\
\addlinespace
& $\text{Mean Deviation}(N)_t = \frac{1}{N}\sum_{i=1}^{N} |\text{TP}_i - \text{SMA}_\text{TP}(N)_t|$ & \\
\addlinespace
Moving Average Convergence Divergence & $\text{MACD}_t = \text{EMA}(N_{\text{short}})_t - \text{EMA}(N_{\text{long}})_t $ & - \\
\addlinespace
& where typically $N_{\text{short}} = 12$ and $N_{\text{long}} = 26$ & \\
\addlinespace
Stochastic RSI & $\text{StochRSI}(N)_t = \frac{\text{RSI}(N)_t - \min(\text{RSI}(N))}{\max(\text{RSI}(N)) - \min(\text{RSI}(N))}$ & [0, 1] \\
\addlinespace
Stochastic Oscillator & $\text{Stoch \%K}(N)_t = \frac{\text{Close}_t - \text{LL}(N)_t}{\text{HH}(N)_t - \text{LL}(N)_t} \times 100$ & [0, 100] \\
\addlinespace
& $\text{Stoch \%D}(M)_t = \text{SMA}_{\text{Stoch \%K}(N)}(M)_t$ & [0, 100] \\
\addlinespace
& $\text{HH}(N)_t = \max_{i=t-N+1}^{t}(\text{High}_i)$, $\text{LL}(N)_t = \min_{i=t-N+1}^{t}(\text{Low}_i)$ & \\
\addlinespace
Ultimate Oscillator & $\text{UO}_t = \frac{4 \times \text{Avg}_t(N_1) + 2 \times \text{Avg}_t(N_2) + \text{Avg}_t(N_3)}{4+2+1} \times 100$ & [0, 100] \\
\addlinespace
& $\text{Avg}_t(N) = \frac{\sum_{i=0}^{N-1}\text{BP}_{t-i}}{\sum_{i=0}^{N-1}\text{TR}_{t-i}}$ & \\
\addlinespace
& $\text{BP}_t = \text{Close}_t - \min(\text{Close}_{t-1}, \text{Low}_t)$ & \\
\addlinespace
& $\text{TR}_t = \max(\text{Close}_{t-1}, \text{High}_t) - \min(\text{Close}_{t-1}, \text{Low}_t)$ & \\
\addlinespace
& where typically $N_1 = 7$, $N_2 = 14$, and $N_3 = 28$ & \\
\addlinespace
Williams \%R & $\text{Williams \%R}(N)_t = \frac{\text{HH}(N)_t - \text{Close}_t}{\text{HH}(N)_t - \text{LL}(N)_t} \times -100$ & [-100, 0] \\
\bottomrule
\end{tabularx}
\caption{Summary of frequently used Momentum Technical Indicators \cite{jansen_machine_2020}.}
\label{Tables:MomentumIndicators}
\end{table}
\begin{table}[htb!]
\centering
\footnotesize
\begin{tabularx}{\textwidth}{@{}lXl@{}}
\toprule
\textbf{Indicator Name} & \textbf{Formula} & \textbf{Range} \\
\midrule
Chaikin Accumulation/Distribution Line & $\text{AD}_t = \text{AD}_{t-1} + \text{MFV}_t$ & - \\
\addlinespace
& $\text{MFV}_t = \frac{(\text{Close}_t - \text{Low}_t) - (\text{High}_t - \text{Close}_t)}{\text{High}_t - \text{Low}_t} \times \text{Volume}_t$ & \\
\addlinespace
Chaikin Accumulation/Distribution Oscillator & $\text{ADOSC}_t = \text{EMA}_\text{AD}(N_{\text{short}})_t - \text{EMA}_\text{AD}(N_{\text{long}})_t$ & - \\
\addlinespace
& where typically $N_{\text{short}} = 3$ and $N_{\text{long}} = 10$ & \\
\addlinespace
On-Balance Volume & $\text{OBV}_t =
\begin{cases}
\text{OBV}_{t-1} + \text{Volume}_t & \text{if } \text{Price}_t > \text{Price}_{t-1} \\
\text{OBV}_{t-1} - \text{Volume}_t & \text{if } \text{Price}_t < \text{Price}_{t-1} \\
\text{OBV}_{t-1} & \text{otherwise}
\end{cases}$ & - \\
\bottomrule
\end{tabularx}
\caption{Summary of frequently used Volume Technical Indicators \cite{jansen_machine_2020}.}
\label{Tables:VolumeIndicators}
\end{table}
\begin{table}[htb!]
\centering
\footnotesize
\begin{tabularx}{\textwidth}{@{}lXl@{}}
\toprule
\textbf{Indicator Name} & \textbf{Formula} & \textbf{Range} \\
\midrule
Standard Deviation & $\text{SD}(N)_t = \sqrt{\frac{1}{N-1} \sum_{i=1}^{N} (\text{Price}_i - \overline{\text{Price}})^2}$ & - \\
\addlinespace
Realised Volatility & $\text{RV}(N)_t = \sqrt{\sigma^2_t}$, where $\sigma^2_t = \alpha \times r_t^2 + (1 - \alpha) \times \sigma^2_{t-1}$, $r_t = \ln\left(\frac{\text{Price}_t}{\text{Price}_{t-1}}\right)$ and $\alpha = \frac{1}{N}$ & - \\
\addlinespace
Bollinger Bands & $\text{Upper BB}(N, k)_t = \text{SMA}(N)_t + k \times \text{SD}(N)_t$ & - \\
\addlinespace
& $\text{Lower BB}(N, k)_t = \text{SMA}(N)_t - k \times \text{SD}(N)_t$ & \\
\addlinespace
& where typically $ k = 2 $& \\
\addlinespace
Average True Range & $\text{ATR}(N)_t = \frac{(N - 1) \times \text{ATR}(N)_{t-1} + \text{TRANGE}_t}{N}$ & - \\
\addlinespace
& $\text{TRANGE}_t = \max(\text{High}_t - \text{Low}_t, |\text{High}_t - \text{Close}_{t-1}|, |\text{Low}_t - \text{Close}_{t-1}|)$ & \\
\addlinespace
Normalised Average True Range & $\text{NATR}(N)_t = \frac{\text{ATR}(N)_t}{\text{Close}_t} \times 100$ & - \\
\addlinespace
Chaikin Volatility Index & $\text{CVI}(N)_t = \frac{\text{EMA}_{HL}(N)_t - \text{EMA}_{HL}(N)_{t-N}}{\text{EMA}_{HL}(N)_{t-N}} \times 100$ & - \\
\addlinespace
& $\text{EMA}_{HL}(N)_t = \text{EMA}(N)$ applied to $(\text{High}_t - \text{Low}_t)$ & \\
\bottomrule
\end{tabularx}
\caption{Summary of frequently used Volatility Technical Indicators \cite{jansen_machine_2020}.}
\label{Tables:VolatilityIndicators}
\end{table}

The four tables survey indicators in common use rather than the contents of any one system. The framework built for this work implements the moving average family and the crossovers between its members, the Rate of Change, the Moving Average Convergence Divergence, the Average True Range, the Realised Volatility and the Efficiency Ratio. It implements no volume indicators, because the tick volume a Forex broker reports counts quote updates rather than traded size, which limits what those indicators can measure. Ten instances of the implemented set form the observation the agent receives, listed in Section~\ref{Section:Implementation}.

Candlestick patterns take a different approach. Rather than deriving a value through a calculation, they describe recognisable shapes formed by one or more candlesticks that traders associate with a change in behaviour \cite{rhoads_candlestick_2022}. Such patterns can be encoded numerically and supplied to a model in the same way as an indicator \cite{jansen_machine_2020}, although this work does not use them: the observation is built from indicators alone.

\subsubsection{Performance Measurement: Returns, Risk and Ratios}

A trading strategy is judged by what it earns and by the risk taken to earn it, so a single return figure is not enough to compare two strategies. The measures below are the ones reported in Section~\ref{Section:ExperimentalResults}.

The total return over a period compares final and initial account equity $E$:

\begin{equation}
R = \frac{E_{final} - E_{initial}}{E_{initial}}
\end{equation}

Because periods differ in length, returns are annualised before they are compared. With a duration $\Delta t$ in seconds and a year taken as $365$ days,

\begin{equation}
R_{annual} = (1 + R)^{\frac{365 \cdot 86400}{\Delta t}} - 1
\label{Equation:AnnualizedReturn}
\end{equation}

Risk is measured in two ways. The first is volatility, the standard deviation of returns, annualised by the rate at which the strategy trades. The second is drawdown. At any moment the drawdown is the distance between current equity and the highest equity reached so far:

\begin{equation}
D(t) = \frac{E(t) - \max_{s \leq t} E(s)}{\max_{s \leq t} E(s)}
\label{Equation:Drawdown}
\end{equation}

The maximum drawdown is the worst value of $D(t)$ over the period and describes the largest loss an investor would have experienced from a peak. The mean drawdown averages $D(t)$ instead, and describes how far below its peak the account sits on a typical day rather than at its worst moment.

Four ratios combine return and risk. All of them divide an annualised excess return by a measure of risk, differing only in which measure they use. The magnitude of that measure is taken, since a drawdown is negative under the definition above. Writing $R_f$ for the risk-free rate,

\begin{equation}
\text{Ratio} = \frac{R_{annual} - R_f}{\left| \text{risk} \right|}
\label{Equation:RiskRatio}
\end{equation}

The \textbf{Sharpe ratio} \cite{sharpe_sharpe_1994} uses annualised volatility, and therefore penalises upward and downward movement equally. The \textbf{Sortino ratio} \cite{sortino_performance_1994} replaces it with downside volatility, computed from negative returns only, on the argument that an investor does not object to upside variability. The \textbf{Calmar ratio} uses the maximum drawdown, so it measures return against the worst loss suffered. The \textbf{Sterling ratio} uses the mean drawdown instead, which makes it less sensitive to a single bad episode and more representative of the typical experience of holding the strategy. Both Calmar and Sterling appear in the literature with several definitions, and the ones used here follow \cite{bacon_practical_2023}. The mean drawdown is smaller than the maximum by construction, so for a strategy with a positive excess return the Sterling ratio is the larger of the two, and the ordering reverses when the excess return is negative.

\subsubsection{Algorithmic and Quantitative Trading}

Financial markets are increasingly dominated by algorithms. Algorithmic trading uses computer programs to execute trades faster and in greater volume than a human can, by analysing large datasets, interpreting market conditions and acting on defined criteria. It is applied to pre-trade analysis, signal generation and trade execution, and it is what allows institutions to run complex strategies at scale. Quantitative trading is a subfield of it, in which mathematical and statistical models are used to find opportunities in market data through historical analysis and machine learning. The approach relies on systematic rules rather than discretionary judgment, and demands knowledge of both finance and mathematics \cite{chan_quantitative_2021}.

Together they have changed how markets operate, introducing practices such as high-frequency trading and shortening the time in which prices react to information. They have also made machine learning central to the field: neural networks for extracting structure from market data, and reinforcement learning for deciding what to do with it \cite{jansen_machine_2020}. Those two are the subject of the sections that follow.